% !PRINTED 366
\weberpage{366}
\gesperrt{Ist $E'$ die Gruppe der in $O'$ enthaltenen Einheiten, und hat $(E, E')$ einen endlichen Wert, so ist die Klassenzahl $h'$ endlich und durch die Formel (12) bestimmt.}

Ist z. B. $O$ die Gruppe aller Zahlen in $\Omega$ (au\ss er 0), $O'$ die Gruppe der Zahlen mit positiver Norm, so ist, falls es \"uberhaupt Zahlen mit negativer Norm gibt, $(O, O') = 2$. Gibt es dann Einheiten mit negativer Norm, so ist auch $(E, E') = 2$. Haben aber alle Einheiten positive Norm, so ist $(E, E') = 1$, und wir erhalten im ersten Falle $h' = h$, im zweiten Falle $2h' = h$.

\gesperrt{Die \"Aquivalenz nach $O$ ist dann die allgemeine \"Aquivalenz, die nach $O'$ die versch\"arfte.}

Kehren wir zu den Ordnungen $[Q]$ zur\"uck und verstehen unter $O$ die zu $Q$ teilerfremden Zahlen (wenn n\"otig nur die mit positiver Norm), so haben wir $(O, O')$ im vorigen Paragraphen bestimmt.

Ist die Diskriminante $D$ und ihr Stamm $\Delta$ negativ, so gibt es im allgemeinen sowohl in $E$ als in $E'$ nur die zwei Einheiten $\pm 1$. In den beiden Ausnahmef\"allen $\Delta = -3, -4$ enth\"alt $E$ sechs und vier Einheiten, $E'$ aber, wenn $Q > 1$ ist, nur zwei, und es ist also $(E, E') = 1$ im allgemeinen, $(E, E') = 3$ im Falle $\Delta = -3$, $(E, E') = 2$ im Falle $\Delta = -4$. Also haben wir in diesen F\"allen nach \S\,98, (10):
\weq{15}{\lambda h' = Q\varPi\left(1 - \frac{(\Delta, q)}{q}\right)h,}
worin:
\weq{16}{
\begin{aligned}
\lambda &= 2, & &\text{f\"ur } \Delta = -4 \\
\lambda &= 3, & &\text{\glqq} \ \Delta = -3 \\
\lambda &= 1, & &\text{\glqq} \ \Delta < -4
\end{aligned}
}
Dies ist die Beziehung, die zuerst von \gesperrt{Gauss} abgeleitet und sp\"ater von \gesperrt{Dirichlet} auf anderem Wege best\"atigt ist.

Ist z. B. $\Delta \equiv 1 \, (\mathrm{mod}\ 4)$, $Q = 2$, so ist nach \gesperrt{Gaussscher} Bezeichnung $h$ die Klassenzahl f\"ur die Formen zweiter, $h'$ die der Formen erster Art. Es ist
\[
\begin{aligned}
Q\left(1 - \frac{(\Delta, q)}{2}\right) &= 1, & \Delta &\equiv 1 \, (\mathrm{mod}\ 8), \\
&= 3, & \Delta &\equiv 5 \, (\mathrm{mod}\ 8).
\end{aligned}
\]
Also
\[
\begin{aligned}
h' &= h, & \Delta &\equiv 1 \\
&= 3h, & \Delta &\equiv 5
\end{aligned}
\qquad (\mathrm{mod}\ 8),
\]
und nur in dem Ausnahmefalle $\Delta = -3$ kommt $h' = h$.
